\documentclass[aspectratio=169, 11pt]{beamer}

% ── Theme ──
\usetheme{Madrid}
\usecolortheme{seahorse}
\setbeamertemplate{navigation symbols}{}
\setbeamertemplate{footline}{%
  \leavevmode\hbox{%
    \begin{beamercolorbox}[wd=.33\paperwidth,ht=2.5ex,dp=1ex,center]{author in head/foot}%
      \usebeamerfont{author in head/foot}\insertshortauthor
    \end{beamercolorbox}%
    \begin{beamercolorbox}[wd=.34\paperwidth,ht=2.5ex,dp=1ex,center]{title in head/foot}%
      \usebeamerfont{title in head/foot}\insertshorttitle
    \end{beamercolorbox}%
    \begin{beamercolorbox}[wd=.33\paperwidth,ht=2.5ex,dp=1ex,right]{date in head/foot}%
      \usebeamerfont{date in head/foot}\insertdate\hspace*{2em}
      \insertframenumber{} / \inserttotalframenumber\hspace*{2ex}
    \end{beamercolorbox}}%
  \vskip0pt%
}

\usepackage{amsmath,amssymb}
\usepackage{graphicx}
\usepackage{booktabs}
\usepackage{tikz}
\usetikzlibrary{arrows.meta, positioning}
\usepackage{hyperref}
\usepackage{xcolor}

\definecolor{accent}{RGB}{0,119,182}
\definecolor{hlgreen}{RGB}{34,139,34}

% ── Meta ──
\title[CS555 Literature Survey]{Real-Time Multi-UAV Incremental Orthomosaicing\\for Flood Response}
\subtitle{CS555 Mid-Semester Presentation --- Literature Survey}
\author[Jugal Alan]{Jugal Alan \\ {\small Supervisor: Dr.\ Shashi Shekhar Jha}}
\institute[IIT Ropar]{M.Tech.\ Artificial Intelligence \\ Indian Institute of Technology Ropar}
\date{February 2026}

\begin{document}

% ============================================================
\begin{frame}
\titlepage
\end{frame}

% ============================================================
\begin{frame}{Outline}
\tableofcontents
\end{frame}

% ============================================================
\section{Introduction \& Context}
% ============================================================

\begin{frame}{Project Context}
\small
\begin{columns}[T]
\column{0.50\textwidth}
\textbf{Our project:} Build a real-time orthomosaic from UAV imagery for flood response

\vspace{0.3em}
\textbf{Core challenges:}
\begin{itemize}
    \item Water surfaces have \textbf{no visual features}
    \item GPS has 1--3\,m drift (consumer receivers)
    \item Must run in real-time ($>$10 FPS)
    \item Must be deployable via web browser
\end{itemize}

\vspace{0.3em}
\textbf{Our approach:}
\begin{itemize}
    \item Direct georeferencing (GPS/IMU)
    \item Tile-based Laplacian blending
    \item Lightweight SIFT-based pose correction
\end{itemize}

\column{0.46\textwidth}
\textbf{5 papers selected to cover:}
\begin{enumerate}
    \item The system we built upon
    \item Modern feature matching (and why we chose classical)
    \item Neural scene representations (and why they don't fit)
    \item State-of-art SLAM (and why it fails on water)
    \item The closest direct competitor
\end{enumerate}

\vspace{0.3em}
\begin{block}{Year range}
\centering
2016 -- 2022 \\ All recent, all highly cited
\end{block}
\end{columns}
\end{frame}

\begin{frame}{Papers at a Glance}
\begin{table}
\centering
\renewcommand{\arraystretch}{1.4}
\begin{tabular}{clll}
\toprule
\textbf{\#} & \textbf{Paper} & \textbf{Venue / Year} & \textbf{Role in Our Work} \\
\midrule
1 & Map2DFusion (Li et al.) & IROS 2016 & Architecture inspiration \\
2 & SuperGlue (Sarlin et al.) & CVPR 2020 & Modern matching baseline \\
3 & NeRF (Mildenhall et al.) & ECCV 2020 & Why not neural representations \\
4 & ORB-SLAM3 (Campos et al.) & IEEE T-RO 2021 & Why not SLAM \\
5 & SfM Revisited (Schonberger & Frahm) & CVPR 2016 & 3D reconstruction baseline \\
\bottomrule
\end{tabular}
\end{table}
\end{frame}

% ============================================================
\section{Paper 1: Map2DFusion}
% ============================================================

\begin{frame}{Paper 1: Map2DFusion --- Overview (Li et al., IROS 2016)}
\small
\textbf{``Real-time Incremental UAV Image Mosaicing Based on Monocular SLAM''}

\begin{columns}[T]
\column{0.50\textwidth}
\textbf{Problem addressed:}
\begin{itemize}
    \item Existing aerial stitching tools (Pix4D, PhotoScan) are \textbf{offline batch} processors
    \item Takes minutes--hours after the flight
    \item No real-time feedback for operators
\end{itemize}

\vspace{0.3em}
\textbf{Core idea:}
\begin{itemize}
    \item Use monocular SLAM (PTAM or LSD-SLAM) for camera pose estimation
    \item Project each frame onto a flat ground plane via inverse pinhole model
    \item Store mosaic as a grid of \textbf{tiles}, each holding \textbf{Laplacian pyramid} coefficients
    \item Only update overlapping tiles $\Rightarrow$ $O(\text{overlap})$ per frame
\end{itemize}

\column{0.46\textwidth}
\textbf{Technical details:}
\begin{itemize}
    \item Tiles: fixed $256{\times}256$ px, hash-indexed by $(t_x, t_y)$
    \item Pyramid: 3--5 frequency bands per tile
    \item Blending: weighted average at each band level
    \item Display: lazy reconstruction --- compose bands only when rendering
    \item Dynamic allocation: new tiles created on demand
\end{itemize}

\vspace{0.3em}
\begin{block}{Reported Performance}
15--25 FPS on a laptop with GPU assistance
\end{block}
\end{columns}
\end{frame}

\begin{frame}{Paper 1: Map2DFusion --- Pipeline}
\small
\textbf{Full pipeline (6 stages):}
\begin{enumerate}
    \item SLAM thread runs in parallel $\rightarrow$ produces 6-DOF camera pose $(R, \mathbf{t})$ at each keyframe
    \item Back-project image corners to the ground plane ($Z{=}0$) using:
    $$\mathbf{P}_W = \mathbf{C} + \lambda \cdot R \cdot K^{-1} \tilde{p}$$
    \item Compute homography $H$ mapping image pixels to ground coordinates
    \item Warp image via $H$; decompose warped patch into Laplacian bands $L_0, L_1, \ldots, L_k$
    \item For each affected tile: merge new bands into stored pyramid via weighted average
    \item On display request: reconstruct tile from pyramid ($\sim$1\,ms)
\end{enumerate}

\vspace{0.3em}
\begin{exampleblock}{Key insight}
By decomposing into frequency bands, high-frequency details (edges) and low-frequency gradients (sky, water) are blended separately --- avoiding ghosting artifacts.
\end{exampleblock}
\end{frame}

\begin{frame}{Paper 1: Map2DFusion --- Limitations}
\small
\begin{columns}[T]
\column{0.50\textwidth}
\begin{alertblock}{Why it doesn't work for floods}
\begin{itemize}
    \item SLAM \textbf{requires visual texture} --- fails over water, sand, snow
    \item Monocular SLAM suffers from \textbf{scale drift} --- the map gradually shrinks or stretches
    \item No GPS anchoring $\Rightarrow$ relative map, not geo-referenced
    \item Single-UAV only --- no protocol for merging multiple sources
\end{itemize}
\end{alertblock}

\column{0.46\textwidth}
\begin{exampleblock}{What we kept vs.\ changed}
\textbf{Kept:}
\begin{itemize}
    \item Tile engine with hash indexing
    \item Laplacian pyramid blending
    \item Lazy display reconstruction
    \item Dynamic tile allocation
\end{itemize}
\textbf{Changed:}
\begin{itemize}
    \item SLAM $\rightarrow$ GPS/IMU direct georeferencing
    \item Added lightweight pose graph
    \item Added multi-UAV support
\end{itemize}
\end{exampleblock}
\end{columns}
\end{frame}

% ============================================================
\section{Paper 2: SuperGlue}
% ============================================================

\begin{frame}{Paper 2: SuperGlue --- Architecture (Sarlin et al., CVPR 2020)}
\small
\textbf{``SuperGlue: Learning Feature Matching with Graph Neural Networks''}

\begin{columns}[T]
\column{0.50\textwidth}
\textbf{Problem addressed:}
\begin{itemize}
    \item Classical matchers (SIFT+NN, ORB+BF) use \textbf{local} descriptor similarity only
    \item No awareness of global geometric layout
    \item Struggle with large viewpoint / lighting changes
\end{itemize}

\vspace{0.3em}
\textbf{Stage 1 --- SuperPoint} (detector/descriptor):
\begin{itemize}
    \item CNN trained on synthetic shapes + homographic adaptation
    \item Outputs repeatable keypoints + 256-D descriptors
    \item More repeatable than SIFT under scale/rotation
\end{itemize}

\column{0.46\textwidth}
\textbf{Stage 2 --- SuperGlue} (matcher):
\begin{itemize}
    \item Graph Neural Network with alternating \textbf{self-attention} and \textbf{cross-attention} layers
    \item Each keypoint ``looks at'' all other keypoints in both images
    \item Final assignment via \textbf{Sinkhorn optimal transport}
\end{itemize}

\vspace{0.3em}
\textbf{Key innovations:}
\begin{itemize}
    \item \textbf{Dustbin:} Learned ``unmatched'' category --- gracefully rejects outliers
    \item \textbf{End-to-end:} Trained on pose-supervised pairs from ScanNet and MegaDepth
\end{itemize}
\end{columns}
\end{frame}

\begin{frame}{Paper 2: SuperGlue --- Benchmark Results}
\small
\begin{columns}[T]
\column{0.50\textwidth}
\textbf{Quantitative results:}
\begin{table}
\centering
\begin{tabular}{lcc}
\toprule
\textbf{Task} & \textbf{SIFT+NN} & \textbf{SuperGlue} \\
\midrule
HPatches AUC@5 & 52.3\% & 70.1\% \\
ScanNet pose AUC & 16.2\% & 31.8\% \\
Outdoor pose AUC & 39.2\% & 59.6\% \\
\bottomrule
\end{tabular}
\end{table}

\vspace{0.3em}
\begin{itemize}
    \item Handles $>$60$^\circ$ viewpoint change
    \item Robust to lighting variations (day $\rightarrow$ night)
    \item 15--20$\times$ fewer outliers than ratio test
\end{itemize}

\column{0.46\textwidth}
\textbf{Computational cost:}
\begin{itemize}
    \item SuperGlue: $\sim$80--120\,ms on RTX 2080 (GPU required)
    \item SIFT + FLANN + ratio test: 20--40\,ms on CPU
\end{itemize}

\vspace{0.3em}
\textbf{Attention mechanism:}
\begin{itemize}
    \item Each keypoint embedding updated by attending to all others
    \item Captures global context (e.g.\ ``this point is to the left of the building'')
    \item Sinkhorn: differentiable optimal transport $\rightarrow$ soft assignment, then discretized
\end{itemize}
\end{columns}
\end{frame}

\begin{frame}{Paper 2: SuperGlue --- Relevance to Our Work}
\small
\begin{columns}[T]
\column{0.50\textwidth}
\textbf{Why this paper matters:}
\begin{itemize}
    \item Sets the \textbf{gold standard} for feature matching accuracy
    \item Our pose graph optimizer uses SIFT matching --- SuperGlue is the natural ``upgrade path''
\end{itemize}

\vspace{0.3em}
\begin{block}{Why we chose SIFT instead}
\begin{itemize}
    \item \textbf{No GPU} required (deploy on CPU-only VMs / edge devices)
    \item Aerial images have high overlap (60--80\%) and small viewpoint change ($<$15$^\circ$) $\Rightarrow$ easy case for SIFT
    \item 2--5$\times$ faster without GPU
\end{itemize}
\end{block}

\column{0.46\textwidth}
\begin{exampleblock}{Future upgrade path}
Plug in SuperGlue for:
\begin{itemize}
    \item Urban scenes with heavy shadows / occlusions
    \item Low-overlap emergency flights
    \item If GPU becomes available on-board
\end{itemize}
\end{exampleblock}

\vspace{0.3em}
\begin{alertblock}{Key takeaway}
SuperGlue proves that learned matchers outperform classical ones on hard cases --- but our aerial scenario is an \textit{easy} case where SIFT suffices at lower cost.
\end{alertblock}
\end{columns}
\end{frame}

% ============================================================
\section{Paper 3: NeRF}
% ============================================================

\begin{frame}{Paper 3: NeRF --- Core Idea (Mildenhall et al., ECCV 2020)}
\small
\textbf{``NeRF: Representing Scenes as Neural Radiance Fields for View Synthesis''}

\begin{columns}[T]
\column{0.50\textwidth}
\textbf{Problem addressed:}
\begin{itemize}
    \item Given posed images, synthesize \textbf{photorealistic novel views}
    \item Previous: explicit meshes, point clouds, voxels --- limited resolution
\end{itemize}

\vspace{0.3em}
\textbf{Implicit neural representation:}
\begin{itemize}
    \item An MLP encodes the entire scene:
    $$F_\Theta: (x,y,z,\theta,\phi) \rightarrow (\mathbf{c}, \sigma)$$
    $\mathbf{c}$ = RGB color, $\sigma$ = volume density
    \item \textbf{Volume rendering} along each camera ray:
    $$C(\mathbf{r}) = \int_{t_n}^{t_f} T(t)\,\sigma(t)\,\mathbf{c}(t)\,dt$$
\end{itemize}

\column{0.46\textwidth}
\textbf{Where:}
$$T(t) = \exp\!\left(-\int_{t_n}^{t}\sigma(s)\,ds\right)$$
$T(t)$: transmittance (how much light reaches point $t$)

\vspace{0.3em}
\textbf{Training:}
\begin{itemize}
    \item Minimize photometric loss:
    $$\mathcal{L} = \sum_{\mathbf{r}} \|C(\mathbf{r}) - C^*(\mathbf{r})\|^2$$
    \item \textbf{Result:} PSNR $>$30\,dB on synthetic scenes; photorealistic novel views of real objects
\end{itemize}
\end{columns}
\end{frame}

\begin{frame}{Paper 3: NeRF --- Technical Details}
\small
\begin{columns}[T]
\column{0.50\textwidth}
\textbf{Positional encoding:}
\begin{itemize}
    \item $\gamma(p) = [\sin(2^k\pi p), \cos(2^k\pi p)]_{k=0}^{L-1}$
    \item Lifts $(x,y,z)$ to high-dimensional space
    \item Allows MLP to capture sharp edges and fine details
\end{itemize}

\vspace{0.3em}
\textbf{Hierarchical sampling:}
\begin{itemize}
    \item Coarse network: uniform 64 samples along ray $\rightarrow$ identifies important regions
    \item Fine network: 128 additional samples concentrated at high-density regions
    \item Reduces computation by 4$\times$ vs.\ uniform dense sampling
\end{itemize}

\column{0.46\textwidth}
\textbf{Aerial / large-scale extensions:}
\begin{itemize}
    \item \textbf{Mega-NeRF} (Turki et al., CVPR 2022): partitions scene into spatial cells, each a sub-NeRF; campus-scale; still hours of training
    \item \textbf{Block-NeRF} (Tancik et al., CVPR 2022): city-block scale; days of training on TPU pods
    \item \textbf{3D Gaussian Splatting} (Kerbl et al., 2023): 100$\times$ faster rendering; still 30--60\,min offline training
\end{itemize}
\end{columns}
\end{frame}

\begin{frame}{Paper 3: NeRF --- Why Not for Our System}
\small
\begin{columns}[T]
\column{0.50\textwidth}
\textbf{NeRF vs.\ Our Explicit Tile Map:}
\begin{table}
\centering
\renewcommand{\arraystretch}{1.15}
\begin{tabular}{lcc}
\toprule
& \textbf{NeRF} & \textbf{Ours} \\
\midrule
Representation & MLP weights & Pixel pyramids \\
Setup time & 12--24 hrs & 0 (streaming) \\
Per-frame & N/A (batch) & 100--210\,ms \\
Hardware & A100 GPU & CPU laptop \\
Output & Any viewpoint & Orthomosaic \\
Incremental & Retrain & $O(1)$ update \\
Use case & Offline insp. & \textbf{Live disaster} \\
\bottomrule
\end{tabular}
\end{table}

\column{0.46\textwidth}
\begin{alertblock}{Fundamental mismatch}
NeRF optimizes for \textbf{visual quality} of novel views.\\[0.3em]
We optimize for \textbf{spatial accuracy} of a 2D map delivered \textbf{in real-time}.\\[0.3em]
Different problems $\Rightarrow$ different tools.
\end{alertblock}

\vspace{0.3em}
\begin{exampleblock}{What we learned}
Neural implicit representations are powerful for reconstruction but fundamentally unsuited for incremental, real-time 2D mapping.
\end{exampleblock}
\end{columns}
\end{frame}

% ============================================================
\section{Paper 4: ORB-SLAM3}
% ============================================================

\begin{frame}{Paper 4: ORB-SLAM3 --- Architecture (Campos et al., IEEE T-RO 2021)}
\small
\textbf{``An Accurate Open-Source Library for Visual, Visual-Inertial, and Multi-Map SLAM''}

\begin{columns}[T]
\column{0.50\textwidth}
\textbf{Problem addressed:}
\begin{itemize}
    \item SLAM from camera streams
    \item ORB-SLAM/v2: mono/stereo only --- no IMU, no recovery from tracking loss
\end{itemize}

\vspace{0.3em}
\textbf{3 parallel threads:}
\begin{enumerate}
    \item \textbf{Tracking:} ORB extraction ($\sim$1000 features/frame), match against local map, PnP + RANSAC pose estimation (30\,FPS)
    \item \textbf{Local Mapping:} Keyframe insertion, triangulation, local bundle adjustment
    \item \textbf{Loop Closing:} DBoW2 place recognition, pose graph optimization, drift correction
\end{enumerate}

\column{0.46\textwidth}
\textbf{Key new features in v3:}
\begin{itemize}
    \item \textbf{Visual-Inertial mode:} Tightly-coupled IMU integration via MAP estimation
    \item \textbf{Atlas multi-map:} On tracking loss, create new sub-map; merge later via place recognition
    \item \textbf{Unified interface:} Mono, stereo, RGB-D, with or without IMU --- all in one system
\end{itemize}
\end{columns}
\end{frame}

\begin{frame}{Paper 4: ORB-SLAM3 --- Benchmarks}
\small
\begin{columns}[T]
\column{0.50\textwidth}
\textbf{Benchmark results:}
\begin{itemize}
    \item EuRoC MAV: \textbf{0.035\,m RMSE} (best published)
    \item TUM-VI: 0.12\,m RMSE
    \item KITTI: comparable to stereo-only systems
\end{itemize}

\vspace{0.3em}
\textbf{Strengths (textured environments):}
\begin{itemize}
    \item Highly accurate: sub-5\,cm indoors
    \item Real-time: 30 FPS tracking
    \item Relocalization: recovers from brief occlusions
    \item Multi-map merging: handles long sequences
\end{itemize}

\column{0.46\textwidth}
\textbf{Indoor vs.\ flood aerial comparison:}
\begin{table}
\centering
\renewcommand{\arraystretch}{1.15}
\begin{tabular}{lcc}
\toprule
\textbf{Issue} & \textbf{Indoor} & \textbf{Flood} \\
\midrule
ORB feat./frame & $\sim$1000 & $<$10 \\
Tracking success & $>$99\% & $<$30\% \\
Scale observable? & Yes & Ambiguous \\
Tracking loss & Rare & Frequent \\
Map output & 3D pts & No mosaic \\
\bottomrule
\end{tabular}
\end{table}
\end{columns}
\end{frame}

\begin{frame}{Paper 4: ORB-SLAM3 --- Why SLAM Fails for Floods}
\small
\begin{columns}[T]
\column{0.50\textwidth}
\begin{alertblock}{4 reasons SLAM doesn't work here}
\begin{enumerate}
    \item \textbf{Feature starvation:} Water/mud have no texture $\Rightarrow$ ORB extraction returns near-zero features
    \item \textbf{Scale drift:} Monocular SLAM can't determine absolute distance --- map shrinks/grows
    \item \textbf{Tracking loss:} Creates fragmented sub-maps that may never merge
    \item \textbf{Wrong output:} SLAM produces sparse 3D point clouds, not dense 2D orthomosaics
\end{enumerate}
\end{alertblock}

\column{0.46\textwidth}
\begin{exampleblock}{What we borrowed}
The \textit{concept} of a pose graph:
\begin{itemize}
    \item Nodes = camera poses
    \item Edges = constraints (GPS + visual matches)
\end{itemize}

\vspace{0.3em}
But simplified:
\begin{itemize}
    \item \textbf{GPS as anchors} (not visual loop closures)
    \item \textbf{2D translation only} (not full 6-DOF)
    \item Works even with zero visual features
\end{itemize}
\end{exampleblock}
\end{columns}
\end{frame}

% ============================================================
\section{Paper 5: Fast UAV Stitching}
% ============================================================

\begin{frame}{Paper 5: Fast UAV Stitching --- Pipeline (Zhang et al., Drones 2022)}
\small
\textbf{``Fast UAV Image Stitching Based on Direct Georeferencing''}

\begin{columns}[T]
\column{0.50\textwidth}
\textbf{Problem addressed:}
\begin{itemize}
    \item SfM-based stitching is too slow for real-time UAV mapping
    \item Can we skip feature matching entirely and use GPS/IMU directly?
\end{itemize}

\vspace{0.3em}
\textbf{Their pipeline:}
\begin{enumerate}
    \item Read GPS coordinates + gimbal angles from EXIF metadata
    \item Compute ground footprint:
    $$\text{GSD} = \frac{h \cdot \text{sensor\_size}}{f \cdot \text{image\_size}}$$
    \item Project image onto geo-referenced canvas via affine transform
    \item Blend overlaps with \textbf{inverse-distance weighting}:
    $$I_{\text{out}}(p) = \frac{\sum_i w_i(p) \cdot I_i(p)}{\sum_i w_i(p)}$$
\end{enumerate}

\column{0.46\textwidth}
\textbf{Results:}
\begin{itemize}
    \item \textbf{10--15 FPS} on agricultural datasets
    \item Tested on DJI Phantom 4 Pro imagery
    \item Positional error: 1--3\,m (consumer GPS)
    \item Single-threaded Python implementation
\end{itemize}

\vspace{0.3em}
\textbf{Advantages:}
\begin{itemize}
    \item Simple and fast
    \item Works over featureless terrain
    \item Deterministic --- same input, same output
    \item No matching failures possible
\end{itemize}
\end{columns}
\end{frame}

\begin{frame}{Paper 5: Fast UAV Stitching --- Limitations}
\small
\begin{columns}[T]
\column{0.50\textwidth}
\textbf{Concrete problems we observed:}
\begin{itemize}
    \item \textbf{Seam artifacts:} Distance-weighted averaging creates ghosting at overlaps
    \item \textbf{No frequency separation:} Sharp edges and smooth gradients handled identically $\Rightarrow$ halos around objects
    \item \textbf{Drift accumulation:} On a 500-image flight, last images off by 3--5\,m from ground truth
    \item \textbf{No redundancy:} Single-UAV, single-pass only
    \item \textbf{No web deployment:} Output is a static image file
\end{itemize}

\column{0.46\textwidth}
\begin{alertblock}{Primary limitation}
No drift correction mechanism --- GPS errors accumulate and create visible seams between adjacent images.
\end{alertblock}

\vspace{0.3em}
\textbf{Blending detail:}
\begin{itemize}
    \item $w_i(p) = 1 / d_i(p)$, where $d_i$ = distance from pixel $p$ to image $i$'s center
    \item Images contribute more at their center, less at edges
    \item But misaligned edges still produce ghosting
\end{itemize}
\end{columns}
\end{frame}

\begin{frame}{Paper 5: Fast UAV Stitching --- Our Improvements}
\small
\begin{columns}[T]
\column{0.50\textwidth}
\textbf{Our improvements over Zhang et al.:}
\begin{table}
\centering
\renewcommand{\arraystretch}{1.15}
\begin{tabular}{lcc}
\toprule
\textbf{Feature} & \textbf{Zhang} & \textbf{Ours} \\
\midrule
Blending & Dist.\ wt.\ avg & Laplacian pyr. \\
Drift correction & None & Pose graph \\

\section{Paper 5: SfM Revisited}
% ============================================================

\begin{frame}{Paper 5: SfM Revisited --- Overview (Schonberger & Frahm, CVPR 2016)}
\small
	extbf{``Structure-from-Motion Revisited''}

\begin{columns}[T]
\column{0.50\textwidth}
	extbf{Problem addressed:}
\begin{itemize}
    \item How to robustly reconstruct 3D structure and camera poses from unordered image collections
    \item Previous SfM pipelines struggled with scalability, robustness, and accuracy
\end{itemize}

\vspace{0.3em}
	extbf{Core contributions:}
\begin{itemize}
    \item Improved feature matching with geometric verification
    \item Robust two-view initialization
    \item Incremental and global bundle adjustment
    \item Efficient outlier filtering and track management
\end{itemize}

\column{0.46\textwidth}
	extbf{Results:}
\begin{itemize}
    \item State-of-the-art accuracy on large-scale datasets
    \item Handles unordered, internet-scale photo collections
    \item Open-source implementation (COLMAP)
\end{itemize}

\vspace{0.3em}
	extbf{Impact:}
\begin{itemize}
    \item Became the standard baseline for 3D reconstruction
    \item Widely used in UAV mapping, photogrammetry, and vision research
\end{itemize}
\end{columns}
\end{frame}

\begin{frame}{Paper 5: SfM Revisited --- Limitations}
\small
\begin{columns}[T]
\column{0.50\textwidth}
	extbf{Limitations in our context:}
\begin{itemize}
    \item Requires sufficient visual features (fails over water, sand, snow)
    \item Computationally intensive (minutes to hours for large datasets)
    \item Not real-time; designed for offline processing
    \item No direct geo-referencing (relative reconstruction)
\end{itemize}

\column{0.46\textwidth}
\begin{alertblock}{Primary limitation}
SfM pipelines cannot handle featureless or dynamic environments, and are not suitable for real-time disaster response mapping.
\end{alertblock}
\end{columns}
\end{frame}

\begin{frame}{Paper 5: SfM Revisited --- Relevance to Our Work}
\small
\begin{columns}[T]
\column{0.50\textwidth}
	extbf{What we learned:}
\begin{itemize}
    \item SfM is the gold standard for offline 3D reconstruction
    \item Our system borrows robust feature matching and outlier rejection ideas
    \item We trade off global accuracy for real-time, incremental updates and water robustness
\end{itemize}

\column{0.46\textwidth}
\begin{exampleblock}{Key takeaway}
SfM is unmatched for high-quality 3D models from images, but is not viable for real-time, water-robust mapping. Our approach adapts some SfM techniques for pose correction, but relies on direct georeferencing for speed and robustness.
\end{exampleblock}
\end{columns}
\end{frame}

